\originalchapter{无碰撞分布函数与 Jeans 质量推断}

本章对应 8.902 第6讲. 知道一个势里的单条轨道还不足以描述星系: 需知道每类轨道上放了多少恒星. 用相空间分布函数记录这种系综, 对其演化方程取速度矩, 才得到观测常用的密度、平均速度与弥散关系. 矩方程简化了问题, 同时留下闭合与各向异性的不确定性.

\originalsection{概率、数密度与质量密度须统一归一化}
本卷用 $f(\boldsymbol x,\boldsymbol v,t)\dd^3x\dd^3v$ 表示该相空间体元中的恒星数. 因而 $\int f\dd^3x\dd^3v=N$, $n(\boldsymbol x)=\int f\dd^3v$. 原课先以随机抽星的概率密度归一化为1, 后又把其速度积分叫数密度; 若使用该概率约定, 应将积分乘 $N$ 才有相同物理数密度.

等质量恒星给 $\rho_*=mn$. 多质量种群需分别定义 $f_a,m_a$, $\rho=\sum_am_a\int f_a\dd^3v$ 再加气体与暗物质. 观测的发光恒星可只是引力势里的示踪种群, 它的 $mn$ 不一定等于生成全部势的 $\rho$.

定义平均速度 $u_i=n^{-1}\int v_if\dd^3v$ 和弥散张量
\[
\sigma_{ij}^2=\overline{v_iv_j}-u_iu_j
=\frac1n\int(v_i-u_i)(v_j-u_j)f\dd^3v.
\]
“零阶矩、一次矩、二次矩”按速度幂次命名, 分别给 $n,nu_i,n\overline{v_iv_j}$. 这样不会把密度误叫含一个速度的第一阶矩.

\originalsection{六维连续性给无碰撞 Boltzmann 方程}
相空间速度为 $(\dot{\boldsymbol x},\dot{\boldsymbol v})=(\boldsymbol v,-\nabla\Phi)$. 粒子数守恒给六维连续方程
\[
\partial_tf+\nabla_x\cdot(f\boldsymbol v)
+\nabla_v\cdot(-f\nabla_x\Phi)=0.
\]
因 $\boldsymbol v$ 不依赖 $\boldsymbol x$, 牛顿势力不依赖 $\boldsymbol v$, 相空间流散度为零, 得
\[
\partial_tf+\boldsymbol v\cdot\nabla_xf
-\nabla_x\Phi\cdot\nabla_vf=0.
\]
这就是沿轨道 $\dd f/\dd t=0$, 与前章相混合的精细守恒相符. 自洽问题还要同求 $\nabla^2\Phi=4\pi G\rho[f]$, 称作 Vlasov--Poisson 系统. 对固定背景中的示踪种群, 则 $\Phi$ 是给定输入. 两者的自由度与稳定性问题不同.

在定常势中, 若 $I_a$ 是轨道积分, 任意非负 $f=F(I_a)$ 都满足定常无碰撞方程, 因为沿轨道每个 $I_a$ 不变. 这给建立平衡模型的充分方法. 球对称非旋转系常取 $f(E,\ell^2)$, 各向同性可取 $f(E)$. 反过来把所有定常解都表示成某个简单有限积分集, 还需轨道的规则性与混合条件; 存在混沌区时不可无条件套最强形式的 Jeans 定理. 定常也不等于稳定, 下一章将显式引入扰动.

\originalsection{逐步取矩得到连续性与动量方程}
假设 $f$ 及 $v_if$ 在速度无穷远衰减足够快, 对无碰撞方程积分. 力项为速度空间全导数, 边界项零, 得
\[
\partial_tn+\partial_i(nu_i)=0.
\]
再乘 $v_j$ 积分. 分部积分给 $\int v_j\partial_{v_i}f\dd^3v=-\delta_{ij}n$, 故
\[
\partial_t(nu_j)+\partial_i[n\overline{v_iv_j}]+n\partial_j\Phi=0.
\]
代入 $\overline{v_iv_j}=u_iu_j+\sigma_{ij}^2$, 用连续方程消去乘 $u_j$ 的项,
\[
\partial_tu_j+u_i\partial_iu_j
=-\frac1n\partial_i(n\sigma_{ij}^2)-\partial_j\Phi.
\]
左侧是平均流的物质导数, 其中对流加速度不是黏性系数; 原笔记将它标成 viscosity/shear, 此处区分几何对流与碰撞耗散. $n\sigma_{ij}^2$ 是随机运动的应力, 可各向异性, 不要求有普通气体热平衡温度. 这些关系是精确的低阶矩方程, 但二次矩的演化会依赖三次矩; 仅有它们并没有闭合求出完整 $f$.

\originalsection{球对称稳态的径向平衡}
对定常、球对称且无平均流的示踪种群, 有 $u_r=u_\theta=u_\varphi=0$, 两个切向弥散相同, 记单一切向分量为 $\sigma_t^2$. 球坐标应力散度的径向部分给
\[
\frac{\dd(n\sigma_r^2)}{\dd r}
+\frac{2n(\sigma_r^2-\sigma_t^2)}r
=-n\frac{GM(<r)}{r^2}.
\]
定义 $\beta=1-\sigma_t^2/\sigma_r^2$（等价于 $1-(\sigma_\theta^2+\sigma_\varphi^2)/(2\sigma_r^2)$）, 得
\[
\frac{\dd(n\sigma_r^2)}{\dd r}+
\frac{2\beta n\sigma_r^2}r=-n\frac{GM(<r)}{r^2}.
\]
$\beta=0$ 等向, $\beta\to1$ 近径向, 大负 $\beta$ 为切向占优. $\sigma_r=0$ 的纯圆轨道极限不宜直接除以它来定义有限 $\beta$, 应回到原应力式取极限.

把导数写成对数斜率, 得质量推断式
\[
M(<r)=-\frac{r\sigma_r^2}{G}
\left[\frac{\dd\ln n}{\dd\ln r}
+\frac{\dd\ln\sigma_r^2}{\dd\ln r}+2\beta\right].
\]
它不是观测数值代入即唯一得到质量的定理, 因为真实三维 $n,\sigma_r,\beta$ 不都可直接测量. 且任意给定函数使此式得到 $M>0$ 只是必要一致性检查, 还未保证存在非负完整分布函数.

当 $\beta$ 固定, 可用积分因子 $r^{2\beta}$ 与外侧应力衰减边界得到
\[
n(r)\sigma_r^2(r)=r^{-2\beta}\int_r^\infty
n(s)GM(<s)s^{2\beta-2}\dd s.
\]
不同边界应力会增加齐次项, 因而截断系统须说明边界压力或轨道支持. 后面的薄盘题也要仔细处理锐利边缘, 不能默认为无任何边界效应.

\originalsection{投影为何产生质量—各向异性退化}
观测给投影半径 $R$ 的表面示踪密度 $\Sigma_n(R)$ 与视线弥散. 对球模型沿视线积分,
\[
\Sigma_n(R)=2\int_R^\infty\frac{n(r)r\dd r}{\sqrt{r^2-R^2}}.
\]
在空间点 $\sin\theta=R/r$, 视线速度二次矩为 $\sigma_r^2\cos^2\theta+\sigma_t^2\sin^2\theta=\sigma_r^2(1-\beta R^2/r^2)$. 因而
\[
\Sigma_n(R)\sigma_{\rm los}^2(R)=2\int_R^\infty
\left(1-\beta\frac{R^2}{r^2}\right)
\frac{n(r)\sigma_r^2(r)r\dd r}{\sqrt{r^2-R^2}}.
\]
同一投影速度数据可以由不同质量和轨道各向异性组合近似重现. 这就是退化的物理来源, 不是代数运算不够熟练. 自行、不同示踪种群、高阶速度分布或独立质量信息能提供额外约束, 但应分别验证模型假设.

\begin{example}
取一个半径区间内 $n\propto r^{-\gamma}$、常数 $\sigma_r,\beta$ 及平坦圆速度, 则 $v_c^2=(\gamma-2\beta)\sigma_r^2$. 若 $\gamma=3$, 同一三维径向弥散在等向假设下给 $v_c^2=3\sigma_r^2$, 在 $\beta=1/2$ 下给 $2\sigma_r^2$. 这个例子显示动力学推断依赖轨道分布; 它尚未要求两模型有完全相同的投影视线弥散, 后者要另用上式计算.
\end{example}

\originalsection{一个可检查的各向同性平衡模型}
令 $f=f_0\exp(-E/\sigma^2)$, $E=v^2/2+\Phi$. 它是轨道能量的函数, 故满足定常无碰撞方程. Gaussian 速度积分给 $n=f_0(2\pi\sigma^2)^{3/2}\exp(-\Phi/\sigma^2)$, 三个速度分量均有弥散 $\sigma^2$. 若其自身产生球势, 奇异等温球的一个形式解为
\[
\Phi=2\sigma^2\ln(r/r_0),\qquad
\rho=\frac{\sigma^2}{2\pi Gr^2},\qquad
v_c^2=2\sigma^2.
\]
由 $\nabla^2\ln r=1/r^2$ 和 $\rho\propto\exp(-\Phi/\sigma^2)$ 双向检查可知它自洽. 但中心奇异、无限总质量和无限延伸使它不是一个完整有限星系模型. 名称“等温”描述等向速度弥散常数, 不意味着无碰撞恒星靠两体热接触已经达到理想气体热平衡. 它把势、密度、分布函数和矩方程放在同一个算例里, 同时展示边界条件不能省略.
